\documentclass[11pt]{article}
\usepackage{amsmath,amssymb,booktabs}
\usepackage[margin=2.3cm]{geometry}
\title{Room of Mathematicians, item 3 (P2 accumulation) — Seat: Tao (analytic)\\
\small Attempt: incomplete Gauss sum / stationary phase for $A_+A_-\neq0$. Not reviewed.}
\date{}
\begin{document}\maketitle

\paragraph{Calibration point.} Read P2\_v2/P2\_v2\_note.tex in full and the block-model section of P2/P2\_note.tex.
Reproduced by hand (mpmath, 50 digits, $N=13$, $a=1$):
\[H_e=\sum_{\rho=0}^{6}\frac{c^\rho\zeta^{\rho^2}}{(\zeta;\zeta)_{2\rho}},\quad
c=-2(1-\zeta),\ \zeta=e^{2\pi i/13},\]
giving $H_e=0.05993-1.93604i$, $H_o=1.26625+1.71861i$, $A_+=0.66309-0.10871i$,
$A_-=-0.60316-1.82732i$ — matches the formula in P2\_v2\_note.tex \S1 exactly (independent
re-derivation, own script, not copy of P2v2\_terms.py/P2v2\_renorm.py). Term-spread confirmed at this
small case too: $|w_\rho|$ over $\rho=0..6$ is $1.0,\,2.15,\,0.94,\,0.243,\,0.063,\,0.0275,\,0.059$ —
non-monotone, no single/edge-dominant term, same qualitative Weyl-spread signature reported for
$N=131$ in P2\_v2\_note.tex \S3. Confirms the prior finding; not retried as a fresh check beyond this.

\paragraph{Task.} Attempt the one alternative route P2\_v2\_note.tex flags as untried: treat $H_e$
(equivalently $H_o$) as an \emph{incomplete} quadratic Gauss sum
\[H_e=\sum_{\rho=0}^{M}w_\rho\,\zeta^{\rho^2},\qquad w_\rho=\frac{c^\rho}{(\zeta;\zeta)_{2\rho}},\qquad
M=\tfrac{N-1}2,\]
and ask whether completion-to-full-period or stationary-phase/van-der-Corput machinery gives an
asymptotic or a nontrivial bound.

\section{Structural check: is $w_\rho$ slowly varying?}
Completion and stationary-phase-with-amplitude methods for incomplete exponential sums (Weyl
differencing, van der Corput, Pólya--Vinogradov-style completion) all require the amplitude
multiplying the phase to be \emph{slowly varying}: bounded total variation, or a bounded number of
monotonicity intervals, on the summation range, so that Abel summation / partial-summation against
the phase's own cancellation is profitable. I checked $|w_\rho|$ directly (own script, $N=131,a=10$,
same case as P2\_v2\_note.tex \S3, extended over the full range $\rho\in[0,65]$):
\begin{center}\small
\begin{tabular}{lllllllllllll}\toprule
$\rho$&0&1&5&10&15&20&21&25&30&40&50&60&65\\
$|w_\rho|$&1.00&2.17&0.026&0.149&0.754&0.565&0.561&0.0079&0.0094&0.035&0.00036&0.0104&0.00027\\\bottomrule
\end{tabular}\end{center}
This is \textbf{not} slowly varying: it spans four orders of magnitude and is non-monotone with at
least $4$--$5$ sign changes of its derivative across the range (up at $\rho=1$, down through $\rho=5$,
up again to a local max near $\rho=15$, down again, a small bump near $\rho=40$, then down to near-zero
at $\rho=50$, back up at $\rho=60$). Reason: $w_\rho=c^\rho/(\zeta;\zeta)_{2\rho}$ is a ratio of a
geometric factor $c^\rho$ (smooth) against a $q$-Pochhammer product $\prod_{j=1}^{2\rho}(1-\zeta^j)$
whose partial products pick up near-zero factors $1-\zeta^j$ whenever $j$ is close to a multiple of
$N$ in the "wrap" sense used by the block model — i.e.\ $w_\rho$ itself already encodes a second,
independent oscillatory/near-singular structure in $\rho$, not a tame envelope.

\section{Consequence for completion}
Completing $\sum_{\rho=0}^M w_\rho\zeta^{\rho^2}$ to a sum over a full period (length $N$ or $2N$,
whichever the quadratic phase's true period is) and bounding the tail requires either (a) $w_\rho$
extending periodically/smoothly past $\rho=M$ so the completed sum is itself meaningful, or (b) the
tail $\sum_{\rho>M} w_\rho\zeta^{\rho^2}$ being controllably small. Neither holds:
\begin{itemize}
\item[(a)] $w_\rho$ has no natural period in $\rho$ matching the phase's period in $\rho^2$; the
denominator $(\zeta;\zeta)_{2\rho}$ is defined by a product up to index $2\rho$ and is not $N$- or
$2N$-periodic in $\rho$ (it strictly accumulates factors). There is no off-the-shelf periodic
extension to complete to.
\item[(b)] The tail is not small: P2\_v2\_note.tex's own table (\S3) and the table above show
$|w_\rho|$ recovering to order $1$ multiple times past small $\rho$ (e.g.\ $\rho=15,20,21$ all
$\gtrsim0.5$, comparable to $\rho=0,1$), so truncating/extending changes the sum by $O(1)$-relative
amounts, not a decaying error.
\end{itemize}
So "complete the incomplete sum" is not merely technically awkward here — the object doesn't have
the shape (slowly-decaying or periodic amplitude) that makes completion an approximation to something
computable; completing it would just be computing a different, unrelated sum.

\section{Stationary phase in $\rho$}
The phase is $\phi(\rho)=a\rho^2/N \pmod 1$, $\phi'(\rho)=2a\rho/N$. As a pure quadratic phase this has
a single stationary point at $\rho=0$ (an endpoint of the sum, not interior), so classical
stationary-phase (which extracts an asymptotic from an interior critical point with $\phi''\neq0$)
degenerates to a boundary/endpoint analysis — exactly the regime where the method is weakest and where
the amplitude's smoothness at the endpoint matters most. Here $w_0=1$ is fine, but the amplitude's wild
behavior a few steps in (Section 1) means even a boundary stationary-phase expansion (Euler-Maclaurin/
endpoint corrections in $\rho$) cannot be truncated at low order: the "error" terms involving
derivatives of $w_\rho$ are not small since $w_\rho$ is not smooth on the scale of the sum.

\section{Verdict}
\textbf{TECHNIQUE DOES NOT APPLY, precisely because of the weight.} The phase $\zeta^{\rho^2}$ is
exactly the right shape for incomplete-Gauss-sum/stationary-phase methods on its own, but the
accompanying weight $w_\rho=c^\rho/(\zeta;\zeta)_{2\rho}$ fails the one hypothesis all of these methods
share (slowly-varying/bounded-variation amplitude, or a periodic/decaying tail to complete against):
it is itself a comparably-sized second oscillatory object, non-monotone over 4 orders of magnitude,
with no periodic extension. This is consistent with, and gives a mechanistic reason for, the
term-spread finding already reported in P2\_v2\_note.tex \S3 (Weyl-type spread, no dominant term): the
spread is not an accident of the phase but is manufactured by the weight's own irregular size.

No asymptotic formula and no nontrivial bound was obtained. This is not merely "harder than expected" —
the specific mechanism (amplitude as irregular as the phase) is the standard obstruction that rules out
van der Corput/completion methods in the literature (they require Lemma-type smoothness hypotheses on
the amplitude that are stated, and violated here, explicitly). A route that might still work would
need to treat $w_\rho\zeta^{\rho^2}$ as a single unified $q$-series exponential-phase object (true
saddle-point analysis of $c^\rho\zeta^{\rho^2}/(\zeta;\zeta)_{2\rho}$ as a whole, in the spirit of the
$q$-Pochhammer-times-quadratic-phase saddle already used for the outer sum in P2\_note.tex, as
suggested in P2\_v2\_note.tex's own "what would actually be needed" paragraph) rather than splitting
into "smooth amplitude times bare Gauss phase" — but that is a different, harder technique (a genuine
$q$-series saddle point) not the elementary incomplete-Gauss-sum/stationary-phase toolkit named in this
brief, and is not attempted here.

\paragraph{Files.} \texttt{small\_hand.py}, \texttt{weight\_check.py} in the session scratchpad
(not committed to the repo; independent re-derivation scripts, mpmath 40--50 digits, each under a
second, run without the perl-alarm wrapper since interactive and well under any risk of hanging;
no BFS/enumeration; negligible memory).

\section*{Weakest points of this note}
\begin{itemize}
\item Only two numeric cases were examined ($N=13,a=1$ by hand; $N=131,a=10$ for the weight table),
both already used or closely mirrored in P2\_v2\_note.tex. The claim that $w_\rho$ is never slowly
varying across the whole window $a/N\in(0,0.0804)$ is not established for other $(a,N)$; it is a
single-case structural diagnosis, offered as a mechanism, not a swept verification.
\item The negative verdict is about the \emph{elementary} incomplete-Gauss-sum/van-der-Corput toolkit
specifically (amplitude-times-phase completion/stationary-phase). It does not rule out a genuine
joint $q$-series saddle-point treatment of $w_\rho\zeta^{\rho^2}$ as one phase (Section 5's suggested
harder alternative), which was not attempted and is a distinct research atom.
\item No attempt was made to check whether a different splitting of the summand (e.g.\ absorbing part
of $(\zeta;\zeta)_{2\rho}$'s phase into an effective quadratic-plus-linear phase and only treating the
remaining magnitude as amplitude) could rescue slow variation; only the literal $w_\rho$ vs.\
$\zeta^{\rho^2}$ split named in the brief was tested.
\end{itemize}
\end{document}
